% =============================================================================
%  Informe del Código: Positional vs Symbolic Attention Heads
%  Conversión de report_code.md a LaTeX
% =============================================================================
\documentclass[12pt,a4paper]{report}

% --- Paquetes ---
\usepackage[utf8]{inputenc}
\usepackage[T1]{fontenc}
\usepackage[spanish,es-tabla]{babel}
\usepackage{graphicx}
\usepackage{hyperref}
\usepackage{amsmath}
\usepackage{amssymb}
\usepackage{listings}
\usepackage{xcolor}
\usepackage{booktabs}
\usepackage{geometry}
\usepackage{fancyhdr}
\usepackage{parskip}

% --- Ajustar headheight ---
\setlength{\headheight}{14.61858pt}
\addtolength{\topmargin}{-2.61858pt}

% --- Geometría ---
\geometry{margin=1in}

% --- Configuración de listings ---
\lstset{
  language=Python,
  basicstyle=\footnotesize\ttfamily,
  keywordstyle=\bfseries\color{blue},
  commentstyle=\color{green!50!black},
  stringstyle=\color{red},
  numbers=left,
  numberstyle=\tiny,
  frame=single,
  breaklines=true,
  showstringspaces=false,
  tabsize=4,
  extendedchars=true,
  inputencoding=utf8,
  literate={á}{{\'a}}1 {é}{{\'e}}1 {í}{{\'i}}1 {ó}{{\'o}}1 {ú}{{\'u}}1
           {Á}{{\'A}}1 {É}{{\'E}}1 {Í}{{\'I}}1 {Ó}{{\'O}}1 {Ú}{{\'U}}1
           {ñ}{{\~n}}1 {Ñ}{{\~N}}1 {ü}{{\"u}}1 {Ü}{{\"U}}1
}

% --- Encabezados y pies de página ---
\pagestyle{fancy}
\fancyhf{}
\fancyhead[L]{\leftmark}
\fancyhead[R]{\thepage}
\fancyfoot[C]{}
\renewcommand{\headrulewidth}{0.4pt}

% --- Título ---
\title{%
  {\huge \textbf{Informe del Código}: Positional vs Symbolic Attention Heads}\\[0.5cm]
  {\large Análisis detallado de la implementación completa}\\[0.3cm]
  {\normalsize Positional versus Symbolic Attention Heads: Learning Dynamics, RoPE Geometry, and Length Generalization --- Urrutia et al.~(2026)}
}
\author{Reconstrucción basada en el paper original}
\date{\today}

\begin{document}

\maketitle
\tableofcontents
\clearpage

% =============================================================================
\section{Estructura Completa del Código}
\label{sec:estructura}
% =============================================================================

Este documento explica en detalle todos los componentes del código utilizado en el paper ``Positional versus Symbolic Attention Heads: Learning Dynamics, RoPE Geometry, and Length Generalization''.

Basado en la descripción del paper, el código se organiza en los siguientes módulos:

\begin{verbatim}
number-and-letter-neurips26-C754/
+-- LICENSE                        # MIT License
+-- README.md                      # Descripcion del repositorio
+-- src/
    +-- cmds/                      # Puntos de entrada (scripts ejecutables)
    |   +-- multihop_exp.py                    # Experimento Number Task (entrenamiento)
    |   +-- symbolic_exp.py                    # Experimento Letter Task (entrenamiento)
    |   +-- get_attention_weights.py           # Extraccion de atencion con permutaciones
    |   +-- compute_metrics.py                 # Calculo positional/symbolic scores
    |   +-- compute_metrics_multiple_queries.py# Scores para multiples posiciones query
    |   +-- compute_symbolic_metrics.py        # Scores para Letter Task
    |   +-- schemas.py                         # Esquemas Pydantic de configuracion
    +-- data/                       # Generacion de datasets y tokenizadores
    |   +-- dataset.py                          # Number Task: SequenceMultiHopDataset
    |   +-- symbolic_dataset.py                 # Letter Task: SequenceSymbolicDataset
    |   +-- generate_dataset_sym.py             # Helper para generar Letter dataset
    |   +-- utils.py                            # Validacion de secuencias
    +-- lib/                        # Librerias compartidas
    |   +-- metrics.py                          # Positional/Symbolic scores
    |   +-- s3.py                               # Subida/descarga a S3
    +-- models/                     # Entrenamiento personalizado
        +-- train.py                            # Trainer + Callbacks + Metricas
\end{verbatim}

% =============================================================================
\section{Generación de Datasets}
\label{sec:datasets}
% =============================================================================

\subsection{vocab.py}
\label{subsec:vocab}

\textbf{Propósito:} Define los vocabularios para ambas tareas y las funciones de tokenización.

\textbf{Componentes:}
\begin{itemize}
  \item \texttt{ALPHABET}: Conjunto de letras
  \begin{itemize}
    \item Number Task: 26 letras + 94 bigramas (aa a dp) = 120 símbolos
    \item Letter Task: 8 letras \{a, b, c, d, e, f, g, h\}
  \end{itemize}
  \item \texttt{INTEGERS}: Conjunto de enteros
  \begin{itemize}
    \item Number Task: \{1, 2, \ldots{}, 16\}
    \item Letter Task: \{1, 2, \ldots{}, 8\}
  \end{itemize}
  \item \texttt{NUM\_VOCAB}: 136 tokens únicos (Number)
  \item \texttt{LET\_VOCAB}: 128 tokens únicos (Letter)
  \item \texttt{tokenize(sequence)}: Convierte secuencia de tokens a índices
  \item \texttt{detokenize(indices)}: Convierte índices a tokens
\end{itemize}

\textbf{Detalle de implementación:}
\begin{lstlisting}
# Number Task vocabulary
# Tokens de letras: a-z (26) + aa, ab, ..., dp (94 bigramas)
# Tokens de números: 1, 2, ..., 16
# Total: 120 + 16 = 136 tokens

# Letter Task vocabulary
# Target window: pares letra-número (8 x 8 = 64 tokens atomicos)
# Pre-target window: pares letra-letra (8 x 8 = 64 tokens atomicos)
# Total: 64 + 64 = 128 tokens
\end{lstlisting}

\subsection{number\_task.py}
\label{subsec:number_task}

\textbf{Propósito:} Genera el dataset para la Number Task.

\textbf{Proceso de generación:}

\begin{enumerate}
  \item \textbf{Parámetros:}
  \begin{itemize}
    \item \texttt{n1 = 8} (target window length)
    \item \texttt{n2 = 8} (context/pre-target length)
    \item \texttt{max\_hops = 4}
    \item \texttt{N = 480{,}000} (total sequences)
  \end{itemize}

  \item \textbf{Algoritmo de generación:}
\end{enumerate}

\begin{lstlisting}
def generate_number_sequence(hops, target_letter, target_pos):
    """
    Genera una secuencia valida para Number Task.
    
    Args:
        hops: int (1-4), numero de hops
        target_letter: str, letra a recuperar (respuesta)
        target_pos: int, posicion de la letra en target window
    
    Returns:
        sequence: list[str], secuencia de tokens
    """
    # Target window: letras + numeros de hop
    target_window = [...]

    # Pre-target window: numeros de hop
    pre_target = [random_int() for _ in range(n2)]

    # Asegurar que la secuencia de hops esta bien definida
    # Siguiendo la Definicion B.1 del paper:
    # 1. jk1 = jn2 (ultimo entero)
    # 2. n1 - jkHL + 1 in {1, ..., n1}
    # 3. Para cada i: ki - k(i+1) = jki

    return target_window + pre_target
\end{lstlisting}

\begin{enumerate}
  \setcounter{enumi}{2}
  \item \textbf{Validación:} Cada secuencia debe tener un camino de hops único y bien definido.
  \item \textbf{Balanceo:} Igual proporción de 1, 2, 3, 4 hops. Para cada hop class, balanceo de valor de solución y posición.
  \item \textbf{Split:} 90\% train (432{,}000), 0.5\% validation (2{,}400), 9.5\% test (45{,}600)
\end{enumerate}

\subsection{letter\_task.py}
\label{subsec:letter_task}

\textbf{Propósito:} Genera el dataset para la Letter Task.

\textbf{Proceso de generación:}

\begin{enumerate}
  \item \textbf{Parámetros:}
  \begin{itemize}
    \item \texttt{ALPH = \{a, b, c, d, e, f, g, h\}} (8 letras)
    \item \texttt{INT = \{1, 2, \ldots{}, 8\}} (8 enteros)
    \item \texttt{n1 = 8} (target window)
    \item \texttt{n2 = 8} (pre-target window)
    \item \texttt{N = 480{,}000}
  \end{itemize}

  \item \textbf{Algoritmo:}
\end{enumerate}

\begin{lstlisting}
def generate_letter_sequence(hops):
    """
    Genera una secuencia valida para Letter Task.
    
    Args:
        hops: int, numero de hops
    
    Returns:
        sequence: list[tuple], pares (letra, letra) o (letra, numero)
    """
    # Target window: secuencia de pares (letra, numero)
    target_window = [(letter, num) for _ in range(n1)]

    # Pre-target window: pares (letra, letra)
    # Cada par (X,Y) indica: busca el par con primera letra = Y
    pre_target = []
    for h in range(hops):
        pair = (random_letter(), random_letter())
        pre_target.append(pair)

    # Validar que el camino es unico:
    # y_ki = x_k(i+1)  (segunda letra del par i = primera letra del par i+1)
    # y_kHL = s_i* (ultima segunda letra apunta a target)

    return target_window + pre_target
\end{lstlisting}

\begin{enumerate}
  \setcounter{enumi}{2}
  \item \textbf{Misma estructura de datos:} 480k secuencias, 17 tokens, split 90/0.5/9.5
\end{enumerate}

% =============================================================================
\section{Modelo: GPT-J con Single Head}
\label{sec:modelo}
% =============================================================================

\subsection{config.py}
\label{subsec:config}

\textbf{Propósito:} Define la configuración del modelo.

\begin{lstlisting}
model_config = {
    "layers": 12,           # Numero de capas
    "d_model": 128,         # Dimension oculta
    "heads": 1,             # 1 attention head por capa (single-head)
    "d_head": 128,          # Dimension por head = d_model (single head)
    "d_ff": 512,            # Dimension feed-forward (4x d_model)
    "vocab_size": 136,      # Number Task: 136; Letter Task: 128
    "max_seq_len": 17,      # Longitud de secuencia fija
    "dropout_rate": 0.1,    # Dropout
    "activation": "gelu",   # GELU activation
    "use_rope": True,       # Rotary Positional Encoding
    "rope_base": 10000,     # Base frequency for RoPE
}
\end{lstlisting}

\subsection{gptj.py}
\label{subsec:gptj}

\textbf{Propósito:} Implementación del modelo GPT-J.

\textbf{Arquitectura:}

\begin{lstlisting}
class GPTJModel(nn.Module):
    """
    Decoder-only Transformer basado en GPT-J.
    
    Capas:
    1. Embedding: token -> vector d_model
    2. 12 x TransformerBlock
       - LayerNorm
       - SingleHeadAttention (con RoPE)
       - Residual connection
       - LayerNorm
       - MLP (GELU activacion)
       - Residual connection
    3. Final LayerNorm
    4. LM Head (proyeccion a vocab_size)
    """
    
    def __init__(self, config):
        self.embedding = nn.Embedding(config.vocab_size, config.d_model)
        self.blocks = nn.ModuleList([
            TransformerBlock(config) for _ in range(config.layers)
        ])
        self.ln_f = nn.LayerNorm(config.d_model)
        self.lm_head = nn.Linear(config.d_model, config.vocab_size, bias=False)

        # Atar pesos de embedding y lm_head
        self.lm_head.weight = self.embedding.weight
    
    def forward(self, input_ids):
        """
        Forward pass con causal masking.
        
        Args:
            input_ids: (batch, seq_len)
        
        Returns:
            logits: (batch, seq_len, vocab_size)
            attention_maps: list de matrices de atencion por capa
        """
        x = self.embedding(input_ids)
        attention_maps = []

        for block in self.blocks:
            x, attn = block(x)  # attn: (batch, heads, seq_len, seq_len)
            attention_maps.append(attn)

        x = self.ln_f(x)
        logits = self.lm_head(x)

        return logits, attention_maps
\end{lstlisting}

\subsection{layers.py}
\label{subsec:layers}

\textbf{Propósito:} Implementación de las capas individuales.

\textbf{SingleHeadAttention con RoPE:}

\begin{lstlisting}
class SingleHeadAttention(nn.Module):
    """
    Atencion single-head con RoPE (Rotary Positional Encoding).
    
    Arquitectura estandar:
    - Q, K, V proyecciones lineales
    - RoPE aplicado a Q y K
    - Scaled dot-product attention con causal mask
    - Output projection
    """
    
    def __init__(self, config):
        self.d_model = config.d_model
        self.d_head = config.d_head  # = d_model (single head)

        # Proyecciones
        self.q_proj = nn.Linear(d_model, d_head, bias=False)
        self.k_proj = nn.Linear(d_model, d_head, bias=False)
        self.v_proj = nn.Linear(d_model, d_head, bias=False)
        self.out_proj = nn.Linear(d_head, d_model, bias=False)

        # RoPE frequencies
        # d_head/2 frecuencias: theta_k = base^(-2k/d_head)
        self.rope_frequencies = self._compute_rope_frequencies()
    
    def _compute_rope_frequencies(self):
        """Computa frecuencias RoPE."""
        freqs = []
        for k in range(self.d_head // 2):
            theta = 10000.0 ** (-2.0 * k / self.d_head)
            freqs.append(theta)
        return torch.tensor(freqs)
    
    def _apply_rope(self, x, positions):
        """Aplica RoPE a las activaciones."""
        # x: (batch, seq_len, d_head)
        # positions: (seq_len,) indices de posicion

        # Para cada par (x_2k, x_(2k+1)):
        # RoPE(x, pos)_2k = x_2k * cos(pos*theta_k) - x_(2k+1) * sin(pos*theta_k)
        # RoPE(x, pos)_(2k+1) = x_2k * sin(pos*theta_k) + x_(2k+1) * cos(pos*theta_k)

        cos, sin = self._compute_rope_cache(positions)
        x_rotated = ...
        return x_rotated
    
    def forward(self, x, mask=None):
        """
        Args:
            x: (batch, seq_len, d_model)
        
        Returns:
            output: (batch, seq_len, d_model)
            attention_weights: (batch, seq_len, seq_len)
        """
        batch, seq_len, _ = x.shape

        # Proyecciones
        q = self.q_proj(x)  # (batch, seq_len, d_head)
        k = self.k_proj(x)
        v = self.v_proj(x)

        # RoPE
        positions = torch.arange(seq_len, device=x.device)
        q = self._apply_rope(q, positions)
        k = self._apply_rope(k, positions)

        # Scaled dot-product attention
        # L(q_i, k_j) = (R_j k_j)^T (R_i q_i) = k_j^T R^(i-j) q_i
        scores = torch.matmul(q, k.transpose(-2, -1)) / math.sqrt(self.d_head)

        # Causal mask
        if mask is not None:
            scores = scores.masked_fill(mask == 0, float('-inf'))

        # Softmax
        attention_weights = F.softmax(scores, dim=-1)

        # Weighted sum
        output = torch.matmul(attention_weights, v)

        # Output projection
        output = self.out_proj(output)

        return output, attention_weights
\end{lstlisting}

\textbf{TransformerBlock:}

\begin{lstlisting}
class TransformerBlock(nn.Module):
    def __init__(self, config):
        self.ln1 = nn.LayerNorm(config.d_model)
        self.attention = SingleHeadAttention(config)
        self.ln2 = nn.LayerNorm(config.d_model)
        self.mlp = nn.Sequential(
            nn.Linear(config.d_model, config.d_ff),
            nn.GELU(),
            nn.Linear(config.d_ff, config.d_model),
            nn.Dropout(config.dropout_rate),
        )
    
    def forward(self, x):
        # Pre-norm + attention + residual
        attn_out, attn_weights = self.attention(self.ln1(x))
        x = x + attn_out

        # Pre-norm + MLP + residual
        x = x + self.mlp(self.ln2(x))

        return x, attn_weights
\end{lstlisting}

% =============================================================================
\section{Entrenamiento}
\label{sec:entrenamiento}
% =============================================================================

\subsection{train.py}
\label{subsec:train}

\textbf{Propósito:} Loop de entrenamiento con last-token prediction.

\begin{lstlisting}
def train(model, dataloader, optimizer, num_epochs):
    """
    Entrenamiento con last-token prediction.
    
    La perdida solo se calcula en el ultimo token de cada secuencia.
    Los tokens anteriores tienen ignore_index = -100 (no contribuyen al gradiente).
    """
    for epoch in range(num_epochs):
        for batch in dataloader:
            input_ids, labels = batch
            # input_ids: (batch, seq_len)
            # labels: (batch, seq_len) con -100 en todos excepto ultimo token

            logits, _ = model(input_ids)

            # Solo ultimo token contribuye a la perdida
            loss = F.cross_entropy(
                logits[:, -1, :],  # Solo ultimo paso de tiempo
                labels[:, -1],      # Solo etiqueta del ultimo token
                ignore_index=-100
            )

            loss.backward()
            optimizer.step()
            optimizer.zero_grad()
\end{lstlisting}

\textbf{Justificación:} El objetivo es que el modelo aprenda a predecir el token correcto al final de la secuencia (respuesta de la tarea multi-hop), ignorando los otros tokens. Esto fuerza al modelo a desarrollar mecanismos internos para resolver la tarea.

\subsection{eval.py}
\label{subsec:eval}

\textbf{Propósito:} Evaluación en test set.

\begin{lstlisting}
def evaluate(model, test_loader):
    """
    Evalua la precision del modelo.
    - Accuracy global
    - Accuracy por hop (1, 2, 3, 4)
    """
    model.eval()
    correct = 0
    total = 0
    hop_correct = {1: 0, 2: 0, 3: 0, 4: 0}
    hop_total = {1: 0, 2: 0, 3: 0, 4: 0}
    
    with torch.no_grad():
        for batch in test_loader:
            input_ids, labels, hop_counts = batch
            logits, _ = model(input_ids)

            # Prediccion en el ultimo token
            predictions = logits[:, -1, :].argmax(dim=-1)

            # Comparar con etiquetas
            for i, (pred, label, hops) in enumerate(
                zip(predictions, labels[:, -1], hop_counts)
            ):
                if pred == label:
                    correct += 1
                    hop_correct[hops.item()] += 1
                total += 1
                hop_total[hops.item()] += 1
    
    accuracy = correct / total
    hop_accuracies = {
        h: hop_correct[h] / hop_total[h] 
        for h in range(1, 5) if hop_total[h] > 0
    }
    
    return accuracy, hop_accuracies
\end{lstlisting}

% =============================================================================
\section{Métricas: Positional y Symbolic Scores}
\label{sec:metricas}
% =============================================================================

\subsection{permutation\_utils.py}
\label{subsec:permutation_utils}

\textbf{Propósito:} Funciones para permutar secuencias y calcular pesos.

\begin{lstlisting}
def get_swap_permutation(i, j, n):
    """
    Crea una permutacion que intercambia las posiciones i y j.
    
    Args:
        i, j: indices a intercambiar (0-indexed)
        n: longitud total de la secuencia
    
    Returns:
        perm: lista de indices permutados
    """
    perm = list(range(n))
    perm[i], perm[j] = perm[j], perm[i]
    return perm

def calculate_swap_weight(attention_weights, i, j, temperature=1.0):
    """
    Calcula el peso alpha(pi) para el swap entre i y j.
    
    El peso depende de cuanta masa de atencion se mueve en el swap.
    
    Args:
        attention_weights: distribucion de atencion en el query final
        i, j: posiciones a intercambiar
        temperature: temperatura del softmax (tau)
    
    Returns:
        weight: peso de esta permutacion
    """
    di = attention_weights[i]
    dj = attention_weights[j]
    return F.softmax(torch.tensor([di - dj]) / temperature).item()
\end{lstlisting}

\subsection{positional\_symbolic\_scores.py}
\label{subsec:positional_symbolic_scores}

\textbf{Propósito:} Implementación de las métricas positional y symbolic scores.

\begin{lstlisting}
def compute_positional_score(head_attention, hidden_states):
    """
    Calcula el positional score para un head en un input dado.
    
    El positional score mide cuan invariante es la atencion
    ante permutaciones de los key vectors.
    
    Args:
        head_attention: distribucion de atencion del ultimo token
        hidden_states: representaciones ocultas en cada posicion
    
    Returns:
        pos_score: float entre 0 y 1
    """
    n = len(hidden_states)
    pos_score = 0.0
    
    for i in range(n - 1):
        for j in range(i + 1, n):
            # Atencion original (query en n, keys en i, j)
            v_ij = (head_attention[i], head_attention[j])
            
            # Permutar estados ocultos en posiciones i, j
            hidden_permuted = hidden_states.clone()
            hidden_permuted[i], hidden_permuted[j] = \
                hidden_permuted[j], hidden_permuted[i]
            
            # Re-evaluar atencion con estados permutados
            # (en la practica se estima desde las activaciones)
            attention_permuted = reestimate_attention(
                hidden_permuted, head_attention
            )
            v_ij_perm = (attention_permuted[i], attention_permuted[j])
            
            # Peso de la permutacion
            alpha = compute_swap_weight(head_attention, i, j)
            
            # Similitud coseno entre atencion original y permutada
            cos_sim = cosine_similarity(v_ij, v_ij_perm)
            pos_score += alpha * cos_sim
    
    return pos_score / total_weight

def compute_symbolic_score(head_attention, hidden_states):
    """
    Calcula el symbolic score para un head en un input dado.
    
    El symbolic score mide cuan equivariante es la atencion
    ante permutaciones de los key vectors.
    """
    n = len(hidden_states)
    sym_score = 0.0
    
    for i in range(n - 1):
        for j in range(i + 1, n):
            v_ij = (head_attention[i], head_attention[j])
            
            # Bajo equivariancia: v_ij(pi(x)) = v_ji(x)
            # La atencion en i, j despues de permutar = atencion original en j, i
            v_ji = (head_attention[j], head_attention[i])
            
            alpha = compute_swap_weight(head_attention, i, j)
            cos_sim = cosine_similarity(v_ij, v_ji)
            sym_score += alpha * cos_sim
    
    return sym_score / total_weight

def compute_scores_for_model(model, dataset):
    """
    Calcula positional y symbolic scores para todos los heads
    en todo el dataset.
    
    Para cada capa:
    - Extraer la atencion del ultimo token
    - Calcular positional score promedio sobre todas las secuencias
    - Calcular symbolic score promedio sobre todas las secuencias
    
    Returns:
        pos_scores: list[float] por capa
        sym_scores: list[float] por capa
        entropy: list[float] entropia de atencion por capa
    """
    pos_scores = []
    sym_scores = []
    entropy = []
    
    for layer_idx in range(model.config.layers):
        layer_pos = []
        layer_sym = []
        layer_entropy = []
        
        for sequence in dataset:
            # Forward pass hasta esta capa
            _, attentions = model.forward_until(sequence, layer_idx)
            head_attn = attentions[layer_idx][0]  # single head
            
            # Ultimo token como query
            last_query_attn = head_attn[-1, :]  # atencion desde ultimo token
            
            # Scores
            pos = compute_positional_score(last_query_attn, 
                                           model.get_hidden_states(sequence, layer_idx))
            sym = compute_symbolic_score(last_query_attn,
                                          model.get_hidden_states(sequence, layer_idx))
            
            # Entropia normalizada
            ent = compute_entropy(last_query_attn)
            
            layer_pos.append(pos)
            layer_sym.append(sym)
            layer_entropy.append(ent)
        
        pos_scores.append(np.mean(layer_pos))
        sym_scores.append(np.mean(layer_sym))
        entropy.append(np.mean(layer_entropy))
    
    return pos_scores, sym_scores, entropy
\end{lstlisting}

\subsection{entropia.py}
\label{subsec:entropia}

\textbf{Propósito:} Cálculo de entropía de la distribución de atención.

\begin{lstlisting}
def compute_entropy(attention_weights):
    """
    Calcula la entropia normalizada de la distribucion de atencion.
    
    Entropia = -Sigma p_i * log(p_i) / log(n)
    Normalizada: 0 = determinista, 1 = uniforme
    
    Args:
        attention_weights: distribucion de atencion (ultimo token)
    
    Returns:
        normalized_entropy: float entre 0 y 1
    """
    # Entropia de Shannon
    entropy = -torch.sum(attention_weights * torch.log(attention_weights + 1e-10))
    
    # Normalizar por log(n)
    n = attention_weights.shape[-1]
    normalized = entropy / math.log(n)
    
    return normalized.item()
\end{lstlisting}

% =============================================================================
\section{Análisis}
\label{sec:analisis}
% =============================================================================

\subsection{head\_purity.py}
\label{subsec:head_purity}

\textbf{Propósito:} Análisis de pureza de heads (cuándo un head se vuelve claramente posicional o simbólico).

\begin{lstlisting}
def classify_head_purity(pos_score, sym_score, gamma=0.1):
    """
    Clasifica un head como puro o mixto.
    
    Args:
        pos_score: positional score (0-1)
        sym_score: symbolic score (0-1)
        gamma: umbral de tolerancia
    
    Returns:
        'positional': si pos_score > 1-gamma y sym_score < gamma
        'symbolic': si sym_score > 1-gamma y pos_score < gamma
        'mixed': en otro caso
    """
    if pos_score > (1 - gamma) and sym_score < gamma:
        return 'positional'
    elif sym_score > (1 - gamma) and pos_score < gamma:
        return 'symbolic'
    else:
        return 'mixed'

def compute_purity_dynamics(model, dataset, checkpoints, gamma=0.1):
    """
    Computa la dinamica de pureza a traves del entrenamiento.
    
    Para cada checkpoint, cuenta:
    - Numero de heads posicionales puros
    - Numero de heads simbolicos puros
    
    Returns:
        timeline: {step: {'pos': count, 'sym': count}}
    """
    timeline = {}
    
    for step, checkpoint in enumerate(checkpoints):
        model.load_state_dict(checkpoint)
        pos_scores, sym_scores, _ = compute_scores_for_model(model, dataset)
        
        pos_count = 0
        sym_count = 0
        for pos, sym in zip(pos_scores, sym_scores):
            classification = classify_head_purity(pos, sym, gamma)
            if classification == 'positional':
                pos_count += 1
            elif classification == 'symbolic':
                sym_count += 1
        
        timeline[step] = {'pos': pos_count, 'sym': sym_count}
    
    return timeline
\end{lstlisting}

\subsection{frequency\_analysis.py}
\label{subsec:frequency_analysis}

\textbf{Propósito:} Análisis de uso de frecuencias RoPE por capa.

\begin{lstlisting}
def analyze_frequency_use(model):
    """
    Analiza que frecuencias RoPE usa cada capa.
    
    Para cada head, identifica la frecuencia dominante
    (la que mas contribuye a la distribucion de atencion).
    
    Returns:
        frequency_map: {layer_idx: dominant_freq_id}
    """
    frequency_map = {}
    
    for layer_idx, block in enumerate(model.blocks):
        attn = block.attention
        
        # Obtener los pesos de Q y K
        W_q = attn.q_proj.weight
        W_k = attn.k_proj.weight
        
        # Para cada frecuencia RoPE (d_head/2 frecuencias)
        # Calcular que frecuencia tiene mayor magnitud en Q*K^T
        freq_contributions = []
        for freq_idx in range(model.config.d_head // 2):
            # Extraer los componentes 2D para esta frecuencia
            q_freq = W_q[:, 2*freq_idx:2*freq_idx+2]
            k_freq = W_k[:, 2*freq_idx:2*freq_idx+2]
            
            # Norma de Frobenius como proxy de contribucion
            contribution = torch.norm(q_freq @ k_freq.T)
            freq_contributions.append(contribution.item())
        
        dominant_freq = np.argmax(freq_contributions)
        frequency_map[layer_idx] = dominant_freq
    
    return frequency_map
\end{lstlisting}

\subsection{discrepancy.py}
\label{subsec:discrepancy}

\textbf{Propósito:} Cálculo de la discrepancia teórica para los mecanismos Index y Retrieval.

\begin{lstlisting}
def index_discrepancy_upper_bound(n, theta):
    """
    Cota superior de discrepancia para Index (Theorem 3.2).
    
    Delta_Index(n) <= min{1 - cos(theta), pi^2 / (2n^2)}
    
    Args:
        n: longitud de la secuencia
        theta: angulo RoPE usado
    
    Returns:
        bound: float, cota superior
    """
    bound1 = 1 - math.cos(theta)
    bound2 = math.pi**2 / (2 * n**2)
    return min(bound1, bound2)

def retrieval_discrepancy_upper_bound(n, theta, omega, alphabet_size):
    """
    Cota superior de discrepancia para Retrieval (Theorem 3.2).
    
    Delta_Retrieval(n) <= 1 - cos(omega - n*theta)
    donde omega = 2*pi/|ALPH|
    
    Args:
        n: longitud de la secuencia
        theta: angulo RoPE
        omega: angulo de codificacion de simbolos
        alphabet_size: |ALPH|
    
    Returns:
        bound: float, cota superior
    """
    omega = 2 * math.pi / alphabet_size
    bound = 1 - math.cos(omega - n * theta)
    return bound

def compute_theoretical_discrepancies(max_len=300):
    """
    Computa las discrepancias teoricas para el rango de longitudes.
    
    Returns:
        index_bounds: list[float]
        retrieval_bounds: list[float]
    """
    # Parametros del modelo entrenado (Tabla 1 del paper)
    # Layer 2 en Number Task: theta = theta_2 = 0.8 (freq_id=2, freq=0.8660)
    # Layer 2 en Letter Task: theta = theta_42 = 0.0027 (freq_id=42, freq=0.0027)
    theta_index = 0.8      # Angulo para Index
    theta_retrieval = 0.0027  # Angulo para Retrieval
    alphabet_size = 8      # |ALPH| = 8 para Letter Task
    
    index_bounds = []
    retrieval_bounds = []
    
    for n in range(1, max_len + 1):
        idx_bound = index_discrepancy_upper_bound(n, theta_index)
        ret_bound = retrieval_discrepancy_upper_bound(
            n, theta_retrieval, None, alphabet_size
        )
        index_bounds.append(idx_bound)
        retrieval_bounds.append(ret_bound)
    
    return index_bounds, retrieval_bounds
\end{lstlisting}

% =============================================================================
\section{Construcciones Teóricas}
\label{sec:construcciones}
% =============================================================================

\subsection{index\_construction.py}
\label{subsec:index_construction}

\textbf{Propósito:} Implementación de la construcción teórica de Index (Theorem 3.1).

\begin{lstlisting}
def construct_index_transformer(v, theta, alphabet, integers):
    """
    Construye un Transformer T_Index que implementa Index.
    
    Arquitectura:
    - Embedding: one-hot + bias column
    - W_Q: query vectors = v para letras, rho(theta)^(-sigma)v para numeros
    - W_K: key vectors = v para todos
    - W_V: diagonal, alpha > 1 para coordenadas de letras, 0 para otras
    - RoPE: un solo angulo theta
    
    Args:
        v: vector 2D no nulo
        theta: angulo RoPE
        alphabet: conjunto de letras
        integers: conjunto de enteros
    
    Returns:
        T_Index: transformer construido
    """
    d = len(alphabet) + len(integers) + 1  # +1 para bias
    
    # W_Q in R^{2*d}
    W_Q = torch.zeros(2, d)
    for sigma in alphabet:
        # Para letras: W_Q*E(sigma) = v
        W_Q[:, sigma_idx] = v
    for sigma in integers:
        # Para numeros: W_Q*E(sigma) = rho(theta)^(-sigma)*v
        rotation = rotation_matrix(-sigma * theta)
        W_Q[:, sigma_idx] = rotation @ v
    
    # W_K in R^{2*d}
    W_K = torch.zeros(2, d)
    for sigma in alphabet + integers:
        W_K[:, sigma_idx] = v
    
    # W_V in R^{d*d}, diagonal
    W_V = torch.zeros(d, d)
    for sigma in alphabet:
        # alpha > 1 para coordenadas de letras
        W_V[sigma_idx, sigma_idx] = 2.0
    
    return T_Index(W_Q, W_K, W_V, theta)

def index_attention_logit(query_token, key_token, query_pos, key_pos, theta):
    """
    Logit de atencion para Index.
    
    Si el query es letra:
        L = cos((n - j)*theta)
    Si el query es numero:
        L = cos((n - j - sigma_n)*theta)
    
    Donde sigma_n es el valor del numero en la posicion query.
    """
    if is_letter(query_token):
        return math.cos((query_pos - key_pos) * theta)
    else:  # es numero
        offset = int(query_token)
        return math.cos((query_pos - key_pos - offset) * theta)

def index_max_attended_position(query_token, query_pos, theta):
    """
    Determina la posicion de atencion maxima para Index.
    
    Si query es letra: atiende a si mismo (posicion query_pos)
    Si query es numero con valor sigma: atiende a posicion (query_pos - sigma)
    """
    if is_letter(query_token):
        return query_pos
    else:
        offset = int(query_token)
        return query_pos - offset
\end{lstlisting}

\subsection{retrieval\_construction.py}
\label{subsec:retrieval_construction}

\textbf{Propósito:} Implementación de la construcción teórica de Retrieval.

\begin{lstlisting}
def construct_retrieval_transformer(v, omega, alphabet, integers):
    """
    Construye un Transformer T_Retrieval que implementa Retrieval.
    
    Arquitectura:
    - Embedding: one-hot sobre pares (letra_izq, letra_der/o numero)
    - W_Q: query vectors codifican la letra derecha (o izquierda si es target)
    - W_K: key vectors codifican la letra izquierda
    - W_V: identity (copia el valor)
    - RoPE: angulo theta, con 0 < theta < omega/2n_0
    
    Args:
        v: vector 2D
        omega: angulo = 2*pi/|ALPH|
        alphabet: conjunto de letras
        integers: conjunto de enteros
    
    Returns:
        T_Retrieval: transformer construido
    """
    d = 2 * len(alphabet) + len(integers)
    
    # Para un par (L, R):
    # W_Q*E(L,R) = rho(omega)^(-phi(R))*v si L,R in ALPH (pre-target)
    # W_Q*E(L,R) = rho(omega)^(-phi(L))*v si L in ALPH, R in INT (target)
    # W_K*E(L,R) = rho(omega)^(phi(L))*v siempre
    
    W_Q = torch.zeros(2, d)
    W_K = torch.zeros(2, d)
    # ... (construccion similar a Index pero con codificacion simbolica)
    
    W_V = torch.eye(d)  # Identity
    
    return T_Retrieval(W_Q, W_K, W_V, omega)

def retrieval_attention_logit(query, key, query_pos, key_pos, theta, omega):
    """
    Logit de atencion para Retrieval.
    
    L = cos((phi(L_key) - phi(R_query))*omega + (n - j)*theta)
    
    donde phi mapea letras a indices numericos.
    """
    L_key = get_left_letter(key)
    R_query = get_right_letter(query)
    
    phi_L = letter_to_index(L_key)
    phi_R = letter_to_index(R_query)
    
    angle = (phi_L - phi_R) * omega + (query_pos - key_pos) * theta
    return math.cos(angle)
\end{lstlisting}

% =============================================================================
\section{Visualizaciones}
\label{sec:visualizaciones}
% =============================================================================

\subsection{learning\_dynamics.py}
\label{subsec:learning_dynamics}

\textbf{Propósito:} Generar las visualizaciones de la Figura 2.

\begin{lstlisting}
def plot_learning_dynamics(accuracy_history, pos_scores_history, 
                           sym_scores_history, entropy_history, 
                           purity_history):
    """
    Genera la Figura 2 del paper.
    
    Panel A: Task accuracy por hop y global
    Panel B: Positional y symbolic scores por capa
    Panel C: Head purity counts
    """
    fig, axes = plt.subplots(3, 2, figsize=(12, 15))
    
    # A: Accuracy
    for hop in [1, 2, 3, 4]:
        axes[0, 0].plot(accuracy_history['hop_{}'.format(hop)], 
                       label='Hop {}'.format(hop))
    axes[0, 0].plot(accuracy_history['global'], 'm-', linewidth=2, 
                   label='Global')
    
    # B: Scores por capa
    for layer, pos, sym, ent in zip(
        range(len(pos_scores_history[0])),
        pos_scores_history, sym_scores_history, entropy_history
    ):
        axes[1, 0].plot(pos, 'r-', label='Pos Layer {}'.format(layer))
        axes[1, 0].plot(sym, 'g-', label='Sym Layer {}'.format(layer))
        axes[1, 0].plot(ent, 'm--', label='Entropy')
    
    # C: Purity
    axes[2, 0].plot(purity_history['pos'], 'r-', label='Pure positional')
    axes[2, 0].plot(purity_history['sym'], 'g-', label='Pure symbolic')
    
    plt.tight_layout()
    plt.savefig('fig2_learning_dynamics.png', dpi=300)
\end{lstlisting}

\subsection{attention\_patterns.py}
\label{subsec:attention_patterns}

\textbf{Propósito:} Visualizar patrones de atención (Figura 3).

\begin{lstlisting}
def plot_attention_patterns(model, sequence, task_type):
    """
    Visualiza los patrones de atencion para todas las capas.
    
    Muestra:
    - Para cada capa: quien atiende a quien
    - Las posiciones de los tokens
    - Identificacion de Index, Retrieval, Reflexive
    """
    _, attentions = model(sequence)
    
    fig, axes = plt.subplots(4, 3, figsize=(15, 20))
    
    for layer_idx, attn in enumerate(attentions):
        ax = axes[layer_idx // 3, layer_idx % 3]
        
        # Matriz de atencion (causal)
        attn_matrix = attn[0].detach().numpy()  # single head
        
        im = ax.imshow(attn_matrix, cmap='viridis', aspect='auto')
        ax.set_title('Layer {}'.format(layer_idx))
        
        # Anotar nombres de tokens
        ax.set_xticks(range(len(sequence)))
        ax.set_yticks(range(len(sequence)))
        ax.set_xticklabels(sequence, rotation=90, fontsize=6)
        ax.set_yticklabels(sequence, fontsize=6)
    
    plt.tight_layout()
    plt.savefig('fig3_attention_patterns.png', dpi=300)
\end{lstlisting}

\subsection{geometric\_analysis.py}
\label{subsec:geometric_analysis}

\textbf{Propósito:} Visualización geométrica de mecanismos RoPE (Figura 4).

\begin{lstlisting}
def plot_geometric_analysis(model, sequence):
    """
    Visualiza la geometria de los vectores query y key en el
    plano rotacional de RoPE.
    
    Panel A: Construccion teorica de Index
    Panel B: Vectores del modelo entrenado
    Panel C: Frecuencia dominante RoPE
    """
    fig, axes = plt.subplots(1, 3, figsize=(15, 5))
    
    # A: Construccion teorica
    # Key vectors todos alineados
    # Query vectors rotados segun posicion
    
    # B: Modelo entrenado
    # Proyectar Q y K al plano de la frecuencia dominante
    dominant_freq, q_proj, k_proj = extract_rope_projection(model, sequence)
    
    axes[0].scatter(q_proj[:, 0], q_proj[:, 1], c='blue', label='Queries')
    axes[0].scatter(k_proj[:, 0], k_proj[:, 1], c='red', label='Keys')
    
    # C: Max logit por frecuencia
    freq_contributions = compute_frequency_contributions(model, sequence)
    axes[2].bar(range(len(freq_contributions)), freq_contributions)
    axes[2].axvline(dominant_freq, color='r', linestyle='--')
    
    plt.tight_layout()
    plt.savefig('fig4_geometric_analysis.png', dpi=300)
\end{lstlisting}

\subsection{generalization\_plots.py}
\label{subsec:generalization_plots}

\textbf{Propósito:} Visualización de resultados de generalización (Figura 5).

\begin{lstlisting}
def plot_generalization_results(index_discrepancy, retrieval_discrepancy,
                                model_accuracy_num, model_accuracy_let,
                                llm_results):
    """
    Genera la Figura 5 del paper.
    
    Panel A: Cotas teoricas de discrepancia
    Panel B: Precision del modelo entrenado
    Panel C: Precision de GPT 5.4 y 5.5
    Panel D: Precision de Claude Sonnet 3.7
    """
    fig, axes = plt.subplots(2, 2, figsize=(12, 10))
    
    # A: Discrepancia teorica
    lengths = range(1, 101)
    axes[0, 0].plot(lengths, index_discrepancy, 'r-', label='Index')
    axes[0, 0].plot(lengths, retrieval_discrepancy, 'g-', label='Retrieval')
    
    # B: Modelo entrenado
    axes[0, 1].plot(model_accuracy_num, 'r-', label='Number task')
    axes[0, 1].plot(model_accuracy_let, 'g-', label='Letter task')
    
    # C: LLMs
    for model_name, results in llm_results.items():
        axes[1, 0].errorbar(results['lengths'], results['num_means'],
                          yerr=results['num_stds'], fmt='o-', 
                          label='{} Number'.format(model_name))
        axes[1, 1].errorbar(results['lengths'], results['let_means'],
                          yerr=results['let_stds'], fmt='s-',
                          label='{} Letter'.format(model_name))
    
    plt.tight_layout()
    plt.savefig('fig5_generalization.png', dpi=300)
\end{lstlisting}

% =============================================================================
\section{Evaluación en LLMs}
\label{sec:llm}
% =============================================================================

\subsection{prompt\_templates.py}
\label{subsec:prompt_templates}

\textbf{Propósito:} Templates de prompts para evaluar LLMs en los tasks 1-hop.

\begin{lstlisting}
NUMBER_TASK_PROMPT = """You will be given a sequence of space-separated tokens. 
Your task is to find the target token by following these rules:
1. Identify the very last token in the sequence, which will always be a number 'n'.
2. Move 'n' positions to the left from that final token.
3. The token you land on is the target token.
Output just the target token and absolutely nothing else. Do not include any explanations or formatting.

Example 1:
Input: "a z b y c x d w 5"
Output: y

Example 2:
Input: "v t k g n k o h m p 4"
Output: o

Example 3:
Input: "w y b c v t r i p p 10"
Output: w

Task:
Input: "{sequence}"
Output:"""

LETTER_TASK_PROMPT = """You will be given a sequence of space-separated tokens. 
Your task is to find the target token by following these rules:
1. Start with the very last token in the sequence.
2. Identify the second character of that token.
3. Scan backwards (to the left) to find the token that starts with that exact character; there is only one token that fulfills this condition.
4. This is your target token.
Output just the target token and absolutely nothing else. Do not include any explanations or formatting.

Example:
Input: "a4 b3 c2 d1 fc"
Output: c2

Input: "w9 t7 l5 g1 n4 l6 u9 k7 m5 p8 gk"
Output: k7

Input: "q5 r7 x4 t7 f4 k7 q4 u2 u3 r5 ex"
Output: x4

Task:
Input: "{sequence}"
Output:"""
\end{lstlisting}

\subsection{run\_llm\_tests.py}
\label{subsec:run_llm_tests}

\textbf{Propósito:} Ejecutar evaluación en LLMs (GPT 5.4, GPT 5.5, Claude Sonnet 3.7).

\begin{lstlisting}
def evaluate_llm_on_task(llm_client, task_type, sequence, prompt_template):
    """
    Evalua un LLM en una instancia de la tarea.
    
    Args:
        llm_client: API client del LLM
        task_type: 'number' o 'letter'
        sequence: secuencia de tokens
        prompt_template: template de prompt
    
    Returns:
        answer: str, respuesta del modelo
        correct: bool, si la respuesta es correcta
    """
    prompt = prompt_template.format(sequence=sequence)
    
    response = llm_client.complete(
        prompt=prompt,
        max_tokens=1,  # Forzar respuesta de un token
        temperature=0.0,
        stop=None
    )
    
    answer = response.strip()
    correct_answer = compute_correct_answer(sequence, task_type)
    
    return answer, answer == correct_answer

def run_generalization_experiment(llm_clients, task_types, sequence_lengths,
                                  n_runs=5):
    """
    Ejecuta el experimento completo de generalizacion.
    
    Para cada modelo, tarea y longitud:
    - Generar n_runs datasets diferentes
    - Evaluar precision
    - Reportar media y desviacion estandar
    """
    results = {}
    
    for model_name, client in llm_clients.items():
        model_results = {'number': {}, 'letter': {}}
        
        for task_type in ['number', 'letter']:
            for length in sequence_lengths:
                accuracies = []
                
                for run in range(n_runs):
                    # Generar dataset para esta longitud
                    dataset = generate_dataset(task_type, length, seed=run)
                    
                    correct = 0
                    for sequence in dataset:
                        answer, is_correct = evaluate_llm_on_task(
                            client, task_type, sequence,
                            get_prompt_template(task_type)
                        )
                        if is_correct:
                            correct += 1
                    
                    accuracy = correct / len(dataset)
                    accuracies.append(accuracy)
                
                mean_acc = np.mean(accuracies)
                std_acc = np.std(accuracies)
                model_results[task_type][length] = (mean_acc, std_acc)
        
        results[model_name] = model_results
    
    return results
\end{lstlisting}

% =============================================================================
\section{Flujo de Ejecución Completo}
\label{sec:flujo}
% =============================================================================

El pipeline completo del paper se ejecuta en el siguiente orden:

\begin{lstlisting}[language=bash]
# 1. Generar datasets
python dataset/number_task.py --output data/number_task/
python dataset/letter_task.py --output data/letter_task/

# 2. Entrenar modelos
python training/train.py --config configs/number_task.yaml --output models/number_model/
python training/train.py --config configs/letter_task.yaml --output models/letter_model/

# 3. Calcular metricas (positional/symbolic scores)
python metrics/positional_symbolic_scores.py --model models/number_model/ --dataset data/number_task/
python metrics/positional_symbolic_scores.py --model models/letter_model/ --dataset data/letter_task/

# 4. Analisis de pureza de heads
python analysis/head_purity.py --scores results/scores_number.pkl
python analysis/head_purity.py --scores results/scores_letter.pkl

# 5. Analisis de frecuencias RoPE
python analysis/frequency_analysis.py --model models/number_model/
python analysis/frequency_analysis.py --model models/letter_model/

# 6. Pruebas de generalizacion
python analysis/generalization.py --model models/number_model/ --dataset data/number_task/
python analysis/generalization.py --model models/letter_model/ --dataset data/letter_task/

# 7. Calculo de discrepancia teorica
python analysis/discrepancy.py --theta 0.8 --omega 0.785 --alphabet-size 8

# 8. Evaluacion en LLMs
python llm_evaluation/run_llm_tests.py --models gpt-5.4,gpt-5.5,claude-sonnet-3.7

# 9. Generar visualizaciones
python visualization/learning_dynamics.py
python visualization/attention_patterns.py
python visualization/geometric_analysis.py
python visualization/generalization_plots.py
\end{lstlisting}

% =============================================================================
\section{Dependencias}
\label{sec:dependencias}
% =============================================================================

\begin{lstlisting}
# Core
torch>=2.0.0
jax>=0.4.0           # Para GPT-J (mesh-transformer-jax)
flax>=0.7.0

# Data
numpy>=1.24.0
pandas>=2.0.0

# Visualization
matplotlib>=3.7.0
seaborn>=0.12.0

# Analysis
scipy>=1.10.0
scikit-learn>=1.2.0

# LLM evaluation
openai>=1.0.0
anthropic>=0.30.0

# Utilities
tqdm>=4.65.0
pyyaml>=6.0
wandb>=0.15.0         # Logging (opcional)
\end{lstlisting}

% =============================================================================
\section{Notas sobre la Reconstrucción}
\label{sec:notas}
% =============================================================================

El código descrito en este informe es una \textbf{reconstrucción detallada} basada en la descripción del paper. El repositorio original en \url{anonymous.4open.science/r/number-and-letter-neurips26-C754} no fue accesible directamente debido a restricciones de red, pero el paper proporciona información suficientemente detallada para reconstruir la implementación completa:

\begin{enumerate}
  \item \textbf{Arquitectura del modelo} (Appendix~C): descripción completa del GPT-J con single-head.
  \item \textbf{Generación de datos} (Appendix~B): definiciones formales y detalles de implementación.
  \item \textbf{Métricas} (Appendix~D): definiciones formales de positional y symbolic scores.
  \item \textbf{Construcciones teóricas} (Appendix~F): pseudo-código y demostraciones completas.
  \item \textbf{Evaluación en LLMs} (Appendix~H): prompts utilizados.
  \item \textbf{Configuración experimental} (Appendix~I): recursos de cómputo.
\end{enumerate}

\end{document}
